%% This is the paper template for the 2022 CFSRC Colloquium 

%%%%%%%%%%%%% packages
%% setup the paper
\documentclass[letter,10pt,twocolumn]{article}
% arial / helvetica font and paper margins
\usepackage[scaled=1]{helvet}
\renewcommand\familydefault{\sfdefault} 
\usepackage[left=0.5in,right=0.5in,top=1in,bottom=1in]{geometry} % margins
% set paragraph spacing and indents
\setlength{\parindent}{0em}
\setlength{\parskip}{1em}
% captions 8.5pt font
\usepackage[font=footnotesize]{caption} % 8.5 pt font
\usepackage[utf8]{inputenc} % for non-ASCII characters

% wallpaper for the conference logo at the top
\usepackage{wallpaper}
\addtolength{\wpXoffset}{-2.90in}
\addtolength{\wpYoffset}{4.6in}

% headers and footers
\usepackage{fancyhdr}
\pagestyle{fancy}
\fancyhf{}
\fancyheadoffset{0cm}
\renewcommand{\headrulewidth}{0pt} 
\renewcommand{\footrulewidth}{0pt}
\fancyfoot{}
\fancyfoot[C]{\fontsize{8}{0}\selectfont \thepage} % page number 8 pt font center bottom

% setup section and subsection title format and numbering
\usepackage{titlesec}
\titlespacing*{\section}{0pt}{0pt}{0pt}
\titlespacing*{\subsection}{0pt}{0pt}{0pt}
\titlespacing*{\subsubsection}{0pt}{0pt}{0pt}
\titleformat{\section}
  {\normalfont\fontsize{10}{10}\bfseries}{\thesection}{1em}{}
 \titleformat{\subsection}
  {\normalfont\fontsize{10}{10}\itshape}{\thesubsection}{1em}{}
 \titleformat{\subsubsection}
  {\normalfont\fontsize{10}{10}\itshape}{\thesubsubsection}{1em}{}
\renewcommand{\thesection}{\arabic{section}.}
\renewcommand{\thesubsection}{\arabic{section}.\arabic{subsection}}

%% other packages for the document
% references style
\usepackage[style=ieee]{biblatex}
\addbibresource{references.bib} % name of reference file
% hide links color and outline
\usepackage[hidelinks]{hyperref}

% packages for tables and figures
\usepackage{float} % floats for figures and tables
\usepackage{multirow} % multiple rows in tables
\usepackage{tabularx} % expandable columns in tables
\usepackage{threeparttable} % table with notes
\usepackage{graphicx} % to input figures
\usepackage{amssymb,amsmath,amsthm,enumitem} % math packages


%%%%%%%%%%%%% start document
\begin{document}\thispagestyle{empty} % remove page number from first page only
\ThisCenterWallPaper{0.2}{cfsrc_short_colloquium_final.png}
% start single column for abstract area
\twocolumn[
  \begin{@twocolumnfalse}

% setup the heading at the top
\begin{flushright}
\vspace{10 pt}
{\small \textit{Proceedings of the Cold-Formed Steel Research Consortium Colloquium}}\\*
{\small \textit{17-19 October, 2022 \ 
{(cfsrc.org)} }}\\*
\end{flushright}

\begin{flushright}
\bigskip
%%%%%%%%%%%%%% Insert your title in the line below
{\large \textbf{CFSRC Colloquium 2022 paper template}}\\* % title
\bigskip
%%%%%%%%%%%%%% Insert your authors in the line below
First A. Author\footnote{}, Second B. Author\footnote{}, Third C. Author\footnote{}\\* % authors
\bigskip
\bigskip
\end{flushright}

\setcounter{secnumdepth}{0} % remove numbering from abstract
% write your abstract
\textbf{Abstract} \\

This is the paper template for the Cold-Formed Steel Research Consortium Colloquium. It comes in two formats, LaTeX and Microsoft Word. This template is required for all papers, and the paper must be provided in PDF format. Provide a brief abstract for your paper, aim for an abstract length between 300-400 words.The abstract should include the objective of the paper, the context in which the work is performed, the methods used in the paper, the basic findings of the work, and perhaps a sense of the future work. Follow normal conventions. Try not to use acronyms or references in the abstract; it should be simple plain text. This paper provides some examples of figures, tables, and equations in this template.

\bigskip
\bigskip

  \end{@twocolumnfalse}] % change to 2 columns

% author affiliations
\footnotetext[1]{Graduate Research Assistant, Department of Civil \& Environmental Engineering, University of Engineering, \href{mailto:emailaddress@email.edu}{emailaddress@email.edu}}
\footnotetext[2]{Engineer, Engineering Company,  \href{mailto:steelcompany@email.com}{steelcompany@email.com}}
\footnotetext[3]{Academic, Department of Civil \& Environmental Engineering, University of Engineering,  \href{mailto:ignore.email@email.edu}{ignore.email@email.edu}}

\setcounter{secnumdepth}{3} % add numbering back in

\section{General}
This section discusses the template style for the main body of the paper. It is recommended to aim for an 8 page paper, although no minimum or maximum paper length is required. The main body font is arial / helectiva 10 point font, and 8.5 point font for the figure and table captions. Please adhere to proper academic writing conventions, write clearly, and try not to forget the use of Oxford commas.

This template aims to cover most common scenarios in conference papers, but may not cover all aspects for your specific paper. Please use your best judgement if any modifications are necessary. This template was written with the use of the LaTeX editor \href{https://www.overleaf.com}{\underline{Overleaf}} in mind, which can be accessed free of charge.

\subsection{Page layout}
This template includes the required margins, font styles, and paragraph spacing. Please refrain from modifying the layout.

\subsubsection{Layers}
You may go up to three layers deep in the organization of your paper, using \textit{section}, \textit{subsection}, and \textit{subsubsection} definitions. The formatting and numbering of these sections is included in the template and no further modifications are required.

\subsection{Units}
International system of units (SI) are required, and other unit systems are optional. The suggested format is metric (US imperial optional) such as 2.3 mm (0.091 inch). This includes units in the main body of the paper, tables, and figures.

\section{Figures}
Single column width figures are preferred, as shown in Figure \ref{logo}. Please make sure that font sizes in the figures are no smaller than 8 point. If needed, figures can be full page width using the \textit{figure*} command, as shown in Figure \ref{logo_large}. Use the word ``Figure" to call out a figure in the text.

\begin{figure}
\centering
    \includegraphics[width=\columnwidth]{cfsrc_short_colloquium_final.png}
	\caption{CFSRC colloquium logo}\vspace*{0pt}
	\label{logo}
\end{figure}

% full page width figure
\begin{figure*}
\centering
    \includegraphics[width=\linewidth]{cfsrc_short_colloquium_final.png}
	\caption{CFSRC colloquium logo large}\vspace*{0pt}
	\label{logo_large}
\end{figure*}

\section{Equations}
Include all necessary equations, such as Equation \ref{circle} shown below for a circle: 
\begin{equation}
    \label{circle}
    A=\pi \times r^2
\end{equation}
% Don't add a blank line here or it will add extra spacing after the equation
where $A$ is the area of a circle and $r$ is the radius of a circle. Use the word ``Equation" to call out the equation in the text.

For longer equations which require multilines, you can use the \textit{multline} command and a sample of this is shown in Equation \ref{sample_eq}.
\begin{multline}
\label{sample_eq}
p(x) = 3x^6 + 14x^5y + 590x^4y^2 + 19x^3y^3\\ % double slash where you want to equation to be split
- 12x^2y^4 - 12xy^5 + 2y^6 - a^3b^3
\end{multline}

\section{Tables}
Single column width tables are preferred, as shown in Table \ref{table_small}. Table font size is 8.5 point font, but may be reduced to 8 point font minimum if needed. If the table is too wide to fit in a single column, please use the \textit{table*} command to make the table full page width. Note that this command will place the large table at the top of the next page. Online \href{https://www.tablesgenerator.com/}{\underline{table generators}} are helpful to use.

\begin{table}[]
	\centering
	\footnotesize %this needed to set the font size to 8.5
	%\scriptsize %this needed to set the font size to 8.0
	\caption{Sample table of numbers}
	\label{table_small}
    \begin{tabularx}{\columnwidth}{>{\centering\arraybackslash}X >{\centering\arraybackslash}X >{\centering\arraybackslash}X >{\centering\arraybackslash}X >{\centering\arraybackslash}X} %this centers each column
    \hline
    Column 1 & Column 2 & Column 3 & Column 4* & Column 5  \\ \hline
    5        & 6        & 7        & 4        & 1         \\ 
    1        & 7        & 3        & 8        & 2         \\ 
    3        & 9        & 6        & 1        & 5         \\ 
    2        & 0        & 8        & 4        & 9         \\ \hline
\end{tabularx}
\begin{tablenotes}
    \footnotesize
    \raggedright
    \item[*] *Sample table footnote for Column 4
\end{tablenotes}
\end{table}

%%%% sample large table
%\begin{table*}[]
%	\centering
%	\footnotesize %this needed to set the font size to 8.5
%	\caption{Sample large table of numbers}
%	\label{table_large}
%   \begin{tabular} {cccccccc}
%    \hline
%   Range & Column 1 & Column 2 & Column 3 & Column 4 & Column 5 & Column 6 & Column 7 \\ \hline
%    0-9   & 2        & 8        & 9        & 3        & 5        & 4        & 7        \\ 
%    10-19 & 19       & 12       & 13       & 11       & 16       & 12       & 14       \\ 
%    20-29 & 28       & 26       & 24       & 29       & 21       & 27       & 22       \\ 
%    30-39 & 34       & 37       & 31       & 36       & 32       & 35       & 38       \\ \hline
%    \end{tabular}
%\end{table*}

\section{References}
The IEEE reference style is used in this template. Please make sure to cite all sources, an example is shown for the AISI Specification \cite{AISIS100-162016}.

\section {Submission}
The paper will be submitted in PDF file format, with the file naming convention of `FirstAuthorFamilyName\_PaperID' sent as an email attachment to \href{mailto:cfsrc@jhu.edu}{cfsrc@jhu.edu}. Please also attach one image (thumbnail) from the paper. Your email subject heading should be ``Paper submission `FirstAuthorFamilyName\_PaperID' colloquium 2022" and should include two attachments - the paper in PDF format and an image. Do not submit your paper or image through a cloud storage link. Submissions are due 1 September 2022.


\section{Conclusions}
Provide a concise set of meaningful conclusions. Please refrain from acronyms and references in the conclusions. Strive for simple plain text. For any urgent questions about this template please email \href{mailto:cfsrc@jhu.edu}{\underline{cfsrc@jhu.edu}}.

\section{Acknowledgments}
Please include any acknowledgements here. Feel free to acknowledge any sponsors or other individuals. Remember to place standard disclaimers as required by funding agencies. 

\printbibliography
\end{document}
